% Reconstructed from the compiled lecture-note PDF.
\chapter{The hidden hyperbolic plane}
\section{The determinant-one surface}

The normalized state matrix X satisfies


\begin{displaytext}
X ∈sl2(R), det X = 1, ⟨J, X⟩< 0.
\end{displaytext}

Write a b ! X = . c −a


\begin{displaytext}
Then −a2 −bc = 1. This is a two-dimensional quadratic surface inside the three-dimensional vector \
space sl2(R). The component selected by ⟨J, X⟩= b −c < 0 is naturally a copy of the hyperbolic \
plane.
\end{displaytext}

\begin{displaytext}
Define the upper half-plane \
H = \{z = x + iy ∈C : y > 0\} .
\end{displaytext}

\begin{displaytext}
For z = x + iy, set \
x/y −(x2 + y2)/y ! \
Φ(z) = . \
1/y −x/y
\end{displaytext}

A direct computation gives


\begin{displaytext}
+ y2 + 1 tr Φ(z) = 0, det Φ(z) = 1, ⟨J, Φ(z)⟩= −x2 < 0. \
y
\end{displaytext}

Conversely, every normalized state matrix in the selected component has this form.


\begin{statementbox}{Theorem}
Theorem 7.1 (Upper-half-plane parametrization). The map
\end{statementbox}

\begin{displaytext}
Φ : H −→\{X ∈sl2(R) : det X = 1, ⟨J, X⟩< 0\}
\end{displaytext}

is a diffeomorphism.


\section{Why matrix conjugation becomes a Möbius transformation}

\begin{displaytext}
A matrix \
α β ! \
g = ∈SL2(R) \
γ δ
\end{displaytext}

\begin{displaytext}
acts on H by \
αz + β \
g · z = . \
γz + δ \
The denominator never vanishes in H, and the imaginary part remains positive. The map Φ is \
equivariant: \
gΦ(z)g−1 = Φ(g · z). \
Thus the conjugation motion of X is exactly the usual action of SL2(R) on the Poincaré upper \
half-plane. \
The stabilizer of J = Φ(i) is the rotation group SO2(R). Hence
\end{displaytext}

H ∼= SL2(R)/SO2(R).


This is the geometric meaning of the adjoint orbit of J.


Background interlude


The Poincaré metric is dx2 + dy2 ds2 = . y2


Its geodesics are vertical lines and semicircles orthogonal to the real axis. Möbius transformations by SL2(R) preserve this metric. The source also uses the symplectic form dx ∧dy/y2, but the metric picture is enough for most of the proof.


\section{The star domain is an ideal triangle}

Substitute X = Φ(x + iy) into the three star inequalities. They become


\begin{displaytext}
−1√ < x < √1 x2 + y2 > 1 3 3, 3.
\end{displaytext}

\begin{displaytext}
Therefore the admissible X-states form \
√ √ \
H∗= nx + iy ∈H : −1/ 3 < x < 1/ 3, x2 + y2 > 1/3 o .
\end{displaytext}

\begin{displaytext}
Its boundary consists of two vertical geodesics and a semicircle orthogonal to the real axis. Its three \
vertices lie at the ideal boundary: \
−1√ √1 ∞. 3, 3, \
Thus H∗is an ideal hyperbolic triangle.
\end{displaytext}

y


H∗


i


\begin{displaytext}
√ √ x \
−1/ 3 1/ 3
\end{displaytext}

\begin{figureplaceholder}
Figure 7.1: The star domain H∗in the upper half-plane. The point z = i corresponds to the circle.
\end{figureplaceholder}

\section{The cost in hyperbolic coordinates}

\begin{displaytext}
Because \
+ y2 + 1 ⟨J, Φ(x + iy)⟩= −x2 , \
y \
the area becomes \
3 Z tf x(t)2 + y(t)2 + 1 \
area(K) = dt. \
2 0 y(t) \
The level sets of \
x2 + y2 + 1 \
y \
are hyperbolic circles centered at i. Indeed,
\end{displaytext}

x2 + y2 + 1 cosh dH(x + iy, i) = . 2y


Thus the cost density is 3 cosh dH(z, i).


It is smallest at the circular state z = i and increases as the state moves hyperbolically away from the center.


\begin{insightbox}{Point requiring care}
\end{insightbox}

\begin{displaytext}
A pointwise smaller cost density does not imply a globally better shape. The circle stays at \
z = i but requires terminal time π/3. Other trajectories can reach the rotational endpoint in \
shorter time. The optimization balances path location and travel time.
\end{displaytext}

\begin{displaytext}
The elementary inequality \
x2 + y2 + 1 1 \
≥y + ≥2 \
y y \
gives \
area(K) ≥3tf.
\end{displaytext}

\begin{displaytext}
Since a minimizer has area less than π, \
π \
tf < . \
3
\end{displaytext}

\section{State equations for (x, y)}

Write au bu ! Zu = , cu −au


\begin{displaytext}
where \
u1 −u2 u0 −2u1 −2u2 \
au = √ , bu = , cu = u0. \
3 3 \
Transferring X′ = [Zu, X]/ ⟨Zu, X⟩through Φ gives
\end{displaytext}

\begin{displaytext}
y(2aux + bu −cux2 + cuy2) x′ = , (7.1) \
2aux + bu −cux2 −cuy2 \
2y2(au −cux) y′ = . (7.2) \
2aux + bu −cux2 −cuy2
\end{displaytext}

Although these formulas look complicated, the vertex trajectories are simple generalized circles.


\section{Constant controls}

Fix u. The quantity ⟨X, Zu⟩is constant because


\begin{displaytext}
d ⟨[Zu, X], Zu⟩ \
⟨X, Zu⟩= = 0. \
dt ⟨Zu, X⟩
\end{displaytext}

Hence Zu P0 = ⟨Zu, X(0)⟩


is constant and X(t) = etP0X(0)e−tP0.


Equivariance implies z(t) = etP0 · z(0).


For a vertex control, P0 has real eigenvalues, and the Möbius orbit is a generalized Euclidean circle joining two ideal vertices of H∗.


\begin{displaytext}
Up to cyclic symmetry: \
√ \
• one vertex gives straight lines from −1/ 3 toward ∞; \
√ \
• a second gives straight lines from ∞toward 1/ 3; \
√ √ \
• the third gives circular arcs from 1/ 3 toward −1/ 3.
\end{displaytext}

e2-flow e1-flow


e0-flow


\begin{figureplaceholder}
Figure 7.2: Schematic vertex-control trajectories in the ideal triangle. Exact trajectories are generalized circles determined by the corresponding Möbius one-parameter groups.
\end{figureplaceholder}

\section{Dihedral symmetry}

The equilateral control triangle has the dihedral symmetry group of order 6 after the central sign in SL2(R) is factored out. Rotation by R cyclically permutes


e0, e1, e2


and the three ideal vertices of H∗. Reflection across the imaginary axis reverses time and exchanges two controls. These symmetries greatly reduce case analysis. A fundamental domain for the action on H∗is one of six sectors bounded by three geodesics through i. Later, a related time-reversal symmetry will exchange the inward and outward chattering fixed points.


\section{Smoothed polygon trajectories}

\begin{displaytext}
A bang-bang control cycles through the three vertices. In H∗, the state follows a concatenation of \
generalized-circle arcs. For the known smoothed (6k + 2)-gons, the orbit repeatedly traces the sides \
of a small equilateral curvilinear triangle centered at i. As k →∞, the triangle shrinks to i, and the \
smoothed polygons converge to the circle.
\end{displaytext}

This hyperbolic picture is not required to define the control problem, but it reveals the geometry of its explicit extremals and identifies the circle as the singular center.


\begin{insightbox}{Chapter summary}
\end{insightbox}

The determinant-one state surface is the hyperbolic upper half-plane. The star inequalities cut out an ideal triangle, and the circle is the central point z = i. The area density is 3 cosh dH(z, i). Constant vertex controls move along generalized circles joining ideal vertices, so bang-bang trajectories become curvilinear polygonal paths in the hyperbolic state space.

